\documentclass[../../main.tex]{subfiles}
\begin{document}

\begin{flushright}
\footnotesize\textit{【自動文字起こし・要確認】}
\end{flushright}

{\bf [解]}
$C_2:x^2+y^2=r^2\ (x,y>0)$とする．まず$C_1$，$C_2$が2点で交わる条件から
\begin{align}
\sqrt{2}<r\quad\cdots①
\end{align}
である．2交点$A(r\cos\alpha,r\sin\alpha)$，$B(r\cos\beta,r\sin\beta)$（$0<\alpha<\beta<\pi/2$）とおくと，対称性から
\begin{align}
\beta=\dfrac{\pi}{2}-\alpha\quad\cdots②
\end{align}
となる．次に，$C_1$が$A$を通る条件から
\begin{align}
r\sin\alpha=\dfrac{1}{r\cos\alpha}\quad\cdots③
\end{align}

$A$における接線$\ell$の傾きは$-\dfrac{1}{r^2\cos^2\alpha}$である．$\angle OAC=\pi/6$（$C$は$\ell$と$x$軸の交点）という条件のもとで角の関係を整理し，$\tan\alpha=t$とおくと
\begin{align}
\dfrac{\sqrt{3}}{3}=\dfrac{t+r^2t+1}{r^2-t^2-t}\quad\cdots④
\end{align}
を得る．ここで，③の両辺に$\dfrac{1}{\cos\alpha}$をかけて整理すると
\begin{align}
r^2t=t^2+1\quad\cdots⑤
\end{align}

④，⑤から
\begin{align}
r^2-t^2-t&=\sqrt{3}(t^2+r^2t+1)\\
t^3+\sqrt{3}t^2+(\sqrt{3}r^2+1)t+\sqrt{3}-r^2&=0\\
(t+\sqrt{3})(r^2t-1)+(\sqrt{3}r^2+1)t+\sqrt{3}-r^2&=0\\
r^2(r^2t-1)+(\sqrt{3}r^2-1)t+(\sqrt{3}r^2+1)t-r^2&=0\\
(r^4+2\sqrt{3}r^2)t&=2r^2
\end{align}
であり，
\begin{align}
t=\dfrac{2}{r^2+2\sqrt{3}}
\end{align}

であり，⑤を代入して，
\begin{align}
t^2+2\sqrt{3}t-1=0
\end{align}
\begin{align}
t=-\sqrt{3}\pm2
\end{align}

$t>0$から複号正を採用して，
\begin{align}
t=2-\sqrt{3}=\tan15^\circ，\quad r=2\quad（①を満たす）\quad\cdots\text{＊}
\end{align}

だから，$0<\alpha<\pi/4$から$\alpha=15^\circ$，⑥である．題意の面積$S$として
\begin{align}
S=(\text{扇形}OAB\text{の面積})-(\text{$C_1$と線分}AB\text{に囲まれた部分の面積})\quad\cdots⑦
\end{align}
であり，
\begin{align}
\text{扇形}OAB&=\dfrac{1}{2}r^2\left(\dfrac{\pi}{2}-\dfrac{\pi}{6}\right)=\dfrac{2}{3}\pi，\\
\text{$C_1$と線分}AB\text{に囲まれた部分}&=\int_{\sin\alpha}^{\cos\alpha}\dfrac{1}{x}\,dx=\log\dfrac{1}{t}
\end{align}

を代入して，＊，⑦から
\begin{align}
S=\dfrac{2}{3}\pi-\log\dfrac{1}{2-\sqrt{3}}=\dfrac{2}{3}\pi-\log(2+\sqrt{3})．
\end{align}

\bigskip
\noindent{\bf [解2]}
対称性から，$y=x$の方を$\ell$とする．$A\left(p,\dfrac{1}{p}\right)$とする（$0<p<1$）．$A$での接線の傾きは$-\dfrac{1}{p^2}$であり，$\ell$と$x$軸の交点を$C$として，$\overline{OA}=\overline{AC}$となる．（$\because$ $OA$の傾き符号）$\angle AOC>\pi/4$から題意の条件は$\angle OAC=\pi/6$である．この時$\angle AOC=\dfrac{5}{12}\pi$となる．
\begin{align}
\tan\dfrac{5}{12}\pi=2+\sqrt{3}
\end{align}
だから，
\begin{align}
\dfrac{1}{p^2}=\tan\dfrac{5}{12}\pi=2+\sqrt{3}\quad\therefore\ p=\sqrt{2-\sqrt{3}}
\end{align}

以下，$\alpha=\sqrt{2-\sqrt{3}}$，$\beta=\sqrt{2+\sqrt{3}}$とする．
\begin{align}
S&=\text{（扇形）}-\text{（$C_1$に囲まれた部分）}\\
&=\dfrac{1}{2}(\alpha^2+\beta^2)\cdot\dfrac{\pi}{3}-\int_\alpha^\beta\dfrac{1}{x}\,dx\\
&=\dfrac{2}{3}\pi-\log(2+\sqrt{3})．
\end{align}

\newpage

\end{document}
